\scanpage{24}
unique) en une \deleted{\unsure{homo}}\inserted{\uline{anti}}dérivation d'Algèbre\inserted{grâce \deleted{à la première} \inserted{aux} conditions \deleted{\unsure{$\beta)$}}\inserted{$\alpha,\beta$}}, et cette dérivation est de carré nul grâce à la condition $\gamma$.

% slanted
\margin{dégager un lemme \unsure{circonstancié}\,\ldots.}

3.\quad \deleted{Algèbre de} \uline{Anneau de Lie} \uline{associé à un complexe de De Rham}

Soit $\mathscr T=\underline{\mathrm{Hom}}(\Omega,\underline O_{X_0})$, c'est un ind-module quasi-ch. sur $X_0$, qu'on identifie à un Module qu. coh.

Ceci définit, si $X,Y$ sont des sections de $\mathscr T$ \deleted{\unsure{sur}} sur un $U$
\[
[X,Y]\deleted{\unsure{(\omega)|(U)}}\in\Gamma(U,\mathscr T)\qquad\text{avec}
\]
\[
\left\langle\underset{\substack{\rotatebox{-90}{$\in$}\\\Gamma(U,\Omega)}}{\omega},[X,Y]\right\rangle\deleted{\unsure{(\omega)}}=-\langle\delta\omega,X\wedge Y\rangle\deleted{\unsure{-}}+\Theta_X\langle\omega,Y\rangle\deleted{\unsure{+}}-\Theta_Y\langle\omega,X\rangle
\tag*{$(x)$}
\]
où l'action $X\mapsto\Theta_X$ de $\mathscr T$ sur $\underline O_{X_0}$ \inserted{(par dérivations)} est définie par
\[
\begin{aligned}
\Theta_X(f)&=\langle\delta f,X\rangle,\qquad [\,\text{NB}\quad\boxed{\Theta_{gX}=g\,\Theta_X}\,]\ \text{résulte de rien}\\
&\hspace{6em}(\text{résulte de }\alpha))\ \boxed{\Theta_X(fg)=\Theta_X(f)g+f\Theta_X(g)}
\end{aligned}
\]
On constate \inserted{grâce à $\alpha)$ et $\beta)$} que le deuxième mb de $(x)$ est bien \uline{$\underline O_{X_0}$-linéaire} en $\omega$, donc définit une section $X\wedge Y$.

De plus, on aura, utilisant $\gamma)$
\[
\deleted{\unsure{[X,Y]}}\ \boxed{\Theta_{[X,Y]}=[\Theta_X,\Theta_Y]}\overset{\mathrm{\unsure{dfn|defn|df}}}{=}\Theta_X\Theta_Y\deleted{\unsure{f}}-\Theta_Y\Theta_X
\tag*{$(xx)$}
\]
\deleted{[\unsure{Si on|Sinon} \unsure{veut|vient} l'autre définition}

[Sans $\gamma$, la \deleted{\unsure{valeur des}} différence des deux mbs sur $f$ est $-\langle\delta\delta f,X\wedge Y\rangle$ ; donc si $\Omega$ loc libre, la condition $\delta_1\delta_0=0$ est \uline{nécessaire} pour qu'on ait $(xx)$]. Ceci vérifie aussi la loi de Lie
\margin{\textcircled{\scriptsize trivial}}
\[
\boxed{[X,X]=0},\qquad\boxed{[[X,Y],Z]+[[Y,Z],X]+[[Z,X],Y]=0}
\]

% archive sheet number: 23

